% 第4章第2节（上）：矩阵的表示（矩阵、元素与形状，设计矩阵与贯穿例）、常见矩阵（方阵、零矩阵、对角矩阵、单位矩阵、上三角矩阵、下三角矩阵、对称矩阵、正定矩阵与半正定矩阵、置换矩阵、正交矩阵）、矩阵的运算（加法与数乘、转置、矩阵与向量的乘积、矩阵乘法、逐项乘积、逆矩阵、矩阵幂、迹、Frobenius内积、Frobenius范数、诱导矩阵范数、行列式）。运算性质随各项运算一并讨论。

\section{矩阵}
\label{sec:ls-matrices}

一批样本的特征排成一张数表以后，问题随之而来：对每个样本用同一组权重求加权和，要算许多个同样形式的和，能否一次写出；两次这样的计算接连进行，能否合成一次；一次计算能否撤销，反复进行时怎样书写；一张表把单位正方形变成平行四边形时，面积放大多少；作为计算，它最多把输入放大多少倍。这些问题都以矩阵为对象。本节先规定矩阵的写法，再介绍元素分布有固定模式的几类常见矩阵，然后逐一定义矩阵的运算，最后按运算整理它们满足的性质。

与向量不同，矩阵最重要的运算不是逐项进行的加法，而是由加权求和导出的乘法。乘法的规则与实数乘法有同有异，需要逐条核对。

\subsection{矩阵的表示与常见矩阵}
\label{sec:ls-matrix-representation}

定位向量中的一个数需要一个索引；一张数表的每个数要由所在的行与列共同确定，需要两个索引。

\paragraph{矩阵、元素与形状}
$m$行$d$列的实数表称为矩阵（matrix），记为$\bm A\in\mathbb R^{m\times d}$；第$i$行、第$j$列的元素记为$\bm A_{i,j}$，不会混淆时也写作$a_{i,j}$。$m\times d$称为矩阵的形状，行数写在列数之前。记号$\bm A_{i,:}$表示第$i$行，$\bm A_{:,j}$表示第$j$列，冒号表示该方向上的全部位置；第$j$列是一个$m$维列向量，第$i$行是一个$d$维行向量。例如在
\[
 \bm A=\begin{bmatrix}1&2&3\\4&5&6\end{bmatrix}\in\mathbb R^{2\times3}
\]
中，$\bm A_{2,1}=4$，第一行$\bm A_{1,:}=\left[\begin{smallmatrix}1&2&3\end{smallmatrix}\right]$，第三列$\bm A_{:,3}=(3,6)^\top$。

两个矩阵相等，要求形状相同且每个对应元素相等。只有一行的矩阵就是行向量，只有一列的矩阵就是列向量，所以向量可以看作矩阵的特例。

\paragraph{设计矩阵}
用矩阵记录数据时，通常让每行对应一个样本，每列对应一个特征，这样的矩阵称为设计矩阵。把示例~\ref{ex:ls-nearly-collinear-features}中两个特征的取值向量$\bm c_1,\bm c_2$并排作为两列，得到三个样本、两个特征的设计矩阵
\begin{equation}
\begin{gathered}
\bm X_{\varepsilon}
 =\begin{bmatrix}1&1\\1&1+\varepsilon\\1&1-\varepsilon\end{bmatrix}
 =\begin{bmatrix}\bm c_1&\bm c_2\end{bmatrix}\\
\bm q=\begin{bmatrix}0\\1\\-1\end{bmatrix},
 \quad \bm c_2=\bm c_1+\varepsilon\bm q.
\end{gathered}
\label{eq:linear-algebra-nearly-collinear-design}
\end{equation}
$|\varepsilon|$很小时，两列之差$\bm c_2-\bm c_1=\varepsilon\bm q$很短，两列几乎是同一串数；第~\ref{sec:ls-inner-product-definition}节将用夹角说明它们几乎同向。这个例子将贯穿整个线性空间部分：第~\ref{chap:linear-representation}章讨论这样两列能否充当一组记录方向，后面的谱与方程两章再看它怎样影响谱、求解与误差。

\subsubsection{常见矩阵}
\label{sec:prelim-common-matrices}

有些矩阵的元素分布有固定模式：大部分位置为零，或者主对角线两侧互相对称。这些模式会让后面的运算大为简化。本小节按元素的位置与取值逐类给出定义和例子；下面随各项运算讨论这些模式如何保留或改变；依赖内积的正交矩阵与正定矩阵在内积空间章建立。

\paragraph{方阵}
行数等于列数的矩阵称为方阵（square matrix），$n\times n$方阵也称$n$阶方阵；位置$(i,i)$上的元素组成主对角线。例如$\left[\begin{smallmatrix}1&2\\3&4\end{smallmatrix}\right]$是二阶方阵，主对角元素为$1$和$4$；$2\times3$矩阵不是方阵；这里讨论三角结构和对角矩阵时均限定为方阵。

\paragraph{零矩阵}
全部元素都为$0$的矩阵称为零矩阵（zero matrix），记作$\bm0$。它不要求是方阵，使用时须写明形状，例如
\[
 \bm0_{2\times3}=\begin{bmatrix}0&0&0\\0&0&0\end{bmatrix}.
\]

\paragraph{对角矩阵}
主对角线以外的元素都为$0$的方阵称为对角矩阵（diagonal matrix），记为$\operatorname{diag}(d_1,\ldots,d_n)$，主对角元素本身可以为$0$。例如
\[
\begin{gathered}
\operatorname{diag}(2,0)=\begin{bmatrix}2&0\\0&0\end{bmatrix}\\
\operatorname{diag}(1,-1,3)=\begin{bmatrix}1&0&0\\0&-1&0\\0&0&3\end{bmatrix}.
\end{gathered}
\]

\paragraph{单位矩阵}
主对角元素全为$1$的对角矩阵称为单位矩阵（identity matrix），记为$\bm I_n$，不必强调阶数时写作$\bm I$。它的第$j$列就是标准坐标向量$\bm e_j$。例如
\[
 \bm I_3=\begin{bmatrix}1&0&0\\0&1&0\\0&0&1\end{bmatrix}=\begin{bmatrix}\bm e_1&\bm e_2&\bm e_3\end{bmatrix}.
\]

\paragraph{上三角矩阵}
$i>j$时$\bm A_{i,j}=0$的方阵称为上三角矩阵（upper triangular matrix）：主对角线以下全为零，非零元素只可能出现在主对角线及其上方。例如
\[
 \begin{bmatrix}1&2&3\\0&4&5\\0&0&6\end{bmatrix}.
\]

\paragraph{下三角矩阵}
$i<j$时$\bm A_{i,j}=0$的方阵称为下三角矩阵（lower triangular matrix）：主对角线以上全为零。例如
\[
 \begin{bmatrix}1&0&0\\2&3&0\\4&5&6\end{bmatrix}.
\]
对角矩阵既是上三角矩阵，也是下三角矩阵。

\paragraph{对称矩阵}
满足$\bm A_{i,j}=\bm A_{j,i}$对全部$i,j$成立的方阵称为对称矩阵（symmetric matrix），它关于主对角线镜像对称；定义转置以后，这个条件写作$\bm A^\top=\bm A$。例如
\[
\begin{bmatrix}1&2\\2&3\end{bmatrix},\qquad
\begin{bmatrix}4&-1&0\\-1&2&5\\0&5&1\end{bmatrix}.
\]
对角矩阵都是对称矩阵。

\paragraph{置换矩阵}
把单位矩阵的各列重新排列得到的方阵称为置换矩阵（permutation matrix），它的每行、每列恰有一个$1$，其余为$0$。例如
\[
\begin{bmatrix}0&1\\1&0\end{bmatrix}=\begin{bmatrix}\bm e_2&\bm e_1\end{bmatrix},\qquad
\begin{bmatrix}0&0&1\\1&0&0\\0&1&0\end{bmatrix}=\begin{bmatrix}\bm e_2&\bm e_3&\bm e_1\end{bmatrix}.
\]
第一个矩阵乘在平面向量上会交换两个坐标，这一点在定义矩阵乘法以后验证。

\subsection{矩阵的基本运算}
\label{sec:ls-matrix-operation-properties}
\label{sec:ls-matrix-operations}
\label{sec:ls-multiplication}

矩阵的运算按来源可以分成几类。加法与数乘沿用向量的逐项规则，转置只改变元素的位置。矩阵与向量的乘积一次写出许多加权和，矩阵乘法、逆矩阵与矩阵幂都以它为基础。逐项乘积是另一种“相乘”，它不汇总任何索引。迹、行列式与矩阵范数把矩阵读成一个数，我们在下一节讨论。本节先给出运算方法，再检查它们满足的规则。

\subsubsection{加法与数乘}

同形矩阵$\bm A,\bm B\in\mathbb R^{m\times n}$逐项相加减，数乘把每个元素乘以同一个实数$c$：
\[
 \begin{gathered}
(\bm A\pm\bm B)_{i,j}=a_{i,j}\pm b_{i,j}\\
(c\bm A)_{i,j}=ca_{i,j},
 \end{gathered}
\]
结果形状不变。下面是一组二阶矩阵的算例。
\begin{example}[矩阵的加法、减法与数乘]
\label{ex:ls-matrix-addition-scaling}
给定：$\bm A=\left(\begin{smallmatrix}1&2\\3&4\end{smallmatrix}\right)$，$\bm B=\left(\begin{smallmatrix}0&1\\-1&2\end{smallmatrix}\right)$，$c=3$。
计算：逐项运算，
\[
\begin{gathered}
\bm A+\bm B=\begin{pmatrix}1&3\\2&6\end{pmatrix}\\
\bm A-\bm B=\begin{pmatrix}1&1\\4&2\end{pmatrix}\\
3\bm A=\begin{pmatrix}3&6\\9&12\end{pmatrix}.
\end{gathered}
\]
例如$\bm A+\bm B$的$(2,1)$元为$3+(-1)=2$，$\bm A-\bm B$的$(2,1)$元为$3-(-1)=4$。
结论：三个结果都是$2\times2$矩阵，每个元素只用到输入在同一位置上的元素。
\end{example}
与向量一样，两个矩阵能够相加，要求形状相同、对应位置含义相同；即使元素总数恰好相同，$2\times3$矩阵与$3\times2$矩阵也不能相加。同形矩阵的线性组合按同样方式定义。

\paragraph{加法与数乘的性质}

同形矩阵逐项相加与数乘，向量一节的八条规则在每个位置上照搬过来。对同形矩阵$\bm A,\bm B,\bm C$和实数$a,b$：
\[
\begin{aligned}
 &\bm A+\bm B=\bm B+\bm A, && (\bm A+\bm B)+\bm C=\bm A+(\bm B+\bm C),\\
 &\bm A+\bm0=\bm A, && \bm A+(-\bm A)=\bm0,\\
 &a(b\bm A)=(ab)\bm A, && 1\bm A=\bm A,\\
 &(a+b)\bm A=a\bm A+b\bm A, && a(\bm A+\bm B)=a\bm A+a\bm B.
\end{aligned}
\]
每一条都可以在任一位置$(i,j)$上用实数运算直接验证。


\subsubsection{转置}

转置交换行与列：$\bm A\in\mathbb R^{m\times n}$的元素$a_{i,j}$原样搬到$\bm A^\top\in\mathbb R^{n\times m}$的位置$(j,i)$，
\[
 (\bm A^\top)_{j,i}=\bm A_{i,j}.
\]
形状交换而数值不变，原来的第$i$行成为第$i$列。把列向量看作$d\times1$矩阵，这一定义与上一节向量的转置一致。下面的例子显示行怎样变成列。
\begin{example}[转置一个$2\times3$矩阵]
\label{ex:ls-transpose-two-by-three}
给定：$\bm A=\left[\begin{smallmatrix}1&2&3\\4&5&6\end{smallmatrix}\right]\in\mathbb R^{2\times3}$。
计算：把每个$(i,j)$元搬到$(j,i)$位置，
\[
 \bm A^\top=\begin{bmatrix}1&4\\2&5\\3&6\end{bmatrix}\in\mathbb R^{3\times2}.
\]
结论：$\bm A$的第一行$(1,2,3)$成为$\bm A^\top$的第一列，第二行$(4,5,6)$成为第二列。
\end{example}

\paragraph{转置的性质}
\label{sec:ls-reshaping}

转置改变元素的位置，也与其他矩阵运算有确定的关系。下面分别检查加法、数乘和乘积。

\paragraph{与加法、数乘交换}
\[
 \begin{gathered}
(\bm A+\bm B)^\top=\bm A^\top+\bm B^\top\\
(c\bm A)^\top=c\bm A^\top\\
(\bm A^\top)^\top=\bm A.
 \end{gathered}
\]
转置只搬动元素，逐项运算的结果搬到新位置后不变；转置两次，每个元素回到原来的位置。

\paragraph{乘积的转置}
\[
 (\bm A\bm B)^\top=\bm B^\top\bm A^\top.
\]
乘积转置后因子次序反转。按元素公式，
\[
 \begin{aligned}
 \bigl((\bm A\bm B)^\top\bigr)_{i,j}
 &=(\bm A\bm B)_{j,i}\\
 &=\sum_\ell a_{j,\ell}b_{\ell,i}\\
 &=\sum_\ell(\bm B^\top)_{i,\ell}(\bm A^\top)_{\ell,j}\\
 &=(\bm B^\top\bm A^\top)_{i,j}.
 \end{aligned}
\]
多个因子时依次使用，$(\bm A\bm B\bm C)^\top=\bm C^\top\bm B^\top\bm A^\top$。

\paragraph{内积中的转置}
\[
\langle\bm A\bm x,\bm y\rangle=\langle\bm x,\bm A^\top\bm y\rangle,\qquad
\bm A\in\mathbb R^{m\times d},\ \bm x\in\mathbb R^d,\ \bm y\in\mathbb R^m.
\]
由上一条，$(\bm A\bm x)^\top\bm y=\bm x^\top\bm A^\top\bm y$：在欧氏内积中，作用在左侧向量上的$\bm A$可以改写成作用在右侧向量上的$\bm A^\top$。

\paragraph{$\bm A^\top\bm A$对称}
对任意矩阵$\bm A$，$\bm A^\top\bm A$与$\bm A\bm A^\top$都是对称矩阵，因为$(\bm A^\top\bm A)^\top=\bm A^\top(\bm A^\top)^\top=\bm A^\top\bm A$，另一个同理。下面直接算出一个$\bm A\bm A^\top$。
\begin{example}[计算$\bm A\bm A^\top$]
\label{ex:ls-gram-product-symmetric}
给定：$\bm A=\left[\begin{smallmatrix}1&2&3\\4&5&6\end{smallmatrix}\right]$。
计算：$\bm A\bm A^\top$的$(i,j)$元是$\bm A$的第$i$行与第$j$行的内积，
\[
 \bm A\bm A^\top=\begin{bmatrix}1+4+9&4+10+18\\4+10+18&16+25+36\end{bmatrix}=\begin{bmatrix}14&32\\32&77\end{bmatrix}.
\]
结论：两个非对角元都是两行的内积$32$，结果是对称矩阵。
\end{example}

\subsubsection{矩阵与向量的乘积}
\label{sec:prelim-matrix-records-prediction}

设$n$次观测的输入排成设计矩阵$\bm X\in\mathbb R^{n\times d}$，权重排成$\bm w\in\mathbb R^d$。线性预测对每次观测计算同一种加权和$\sum_j x_{i,j}w_j$，即$\bm X$的第$i$行与$\bm w$的内积。矩阵与向量的乘积把这$n$个加权和一次写出：对$\bm A\in\mathbb R^{m\times d}$与$\bm u\in\mathbb R^d$，$\bm A\bm u$是$m$维列向量，
\begin{equation}
\begin{gathered}
(\bm A\bm u)_i=\sum_{j=1}^d\bm A_{i,j}u_j\\
i=1,\ldots,m.
\end{gathered}
\label{eq:prelim-matrix-prediction}
\end{equation}
矩阵的列数必须等于向量的维数。同一个结果有两种算法：
\begin{itemize}
 \item 按行算：输出的第$i$个坐标是第$i$行与$\bm u$的内积；
 \item 按列算：同一个$u_j$乘的恰好是第$j$列的全部元素，所以
 \[
  \bm A\bm u=u_1\bm A_{:,1}+u_2\bm A_{:,2}+\cdots+u_d\bm A_{:,d},
 \]
 乘积是矩阵各列以$u_j$为系数的线性组合。
\end{itemize}
特别地，按列算时，$\bm e_j$只有第$j$个系数为$1$，所以$\bm A\bm e_j=\bm A_{:,j}$：乘第$j$个标准坐标向量就取出第$j$列。

下面用一个二阶例子核对两种算法给出同一结果。
\begin{example}[按行与按列计算矩阵与向量的乘积]
\label{ex:ls-matrix-vector-rows-columns}
给定：$\bm A=\left[\begin{smallmatrix}2&1\\0&3\end{smallmatrix}\right]$，$\bm u=(1,1)^\top$。
计算：按行算，每个坐标是一行与$\bm u$的内积；按列算，结果是两列以$u_1=u_2=1$为系数的线性组合：
\[
\begin{gathered}
\bm A\bm u
 =\begin{bmatrix}2\cdot1+1\cdot1\\0\cdot1+3\cdot1\end{bmatrix}
 =\begin{bmatrix}3\\3\end{bmatrix}\\
\bm A\bm u=1\begin{bmatrix}2\\0\end{bmatrix}+1\begin{bmatrix}1\\3\end{bmatrix}
 =\begin{bmatrix}3\\3\end{bmatrix}.
\end{gathered}
\]
结论：两种算法都得到$\bm A\bm u=(3,3)^\top$。
\end{example}

用于数据时，矩阵与向量的乘积一次给出全部样本的预测。下面以一次函数为例。
\begin{example}[用矩阵与向量的乘积计算一次函数的取值]
\label{ex:ls-affine-predictions}
给定：要在$x=1,2,3$处计算$ax+b$，取
\[
\bm X=\begin{pmatrix}1&1\\2&1\\3&1\end{pmatrix},\qquad
\bm w=\begin{pmatrix}a\\b\end{pmatrix}.
\]
计算：按行算，
\[
 \bm X\bm w=\begin{pmatrix}1\cdot a+1\cdot b\\2\cdot a+1\cdot b\\3\cdot a+1\cdot b\end{pmatrix}=\begin{pmatrix}a+b\\2a+b\\3a+b\end{pmatrix}.
\]
结论：$\bm X\bm w$的三个坐标正是三处的函数值。第二列全为$1$，使每行都加上同一个截距$b$。
\end{example}

观测输出排成$\bm y$以后，全部残差就是$\bm X\bm w-\bm y$，平方误差总和写成
\[
 \sum_{i=1}^n\left(\sum_{j=1}^d x_{i,j}w_j-y_i\right)^2
 =\lVert\bm X\bm w-\bm y\rVert_2^2.
\]
左侧逐个样本求加权和、减去观测、平方后相加；右侧把同样的计算写成残差向量长度的平方。对贯穿例，按列算得$\bm X_\varepsilon\bm w=w_1\bm c_1+w_2\bm c_2$：预测是两个特征取值向量的线性组合。

\subsubsection{矩阵乘法}
\phantomsection
\label{sec:prelim-linear-actions-composition}

对$\bm A\in\mathbb R^{m\times d}$与$\bm B\in\mathbb R^{d\times k}$，乘积$\bm A\bm B\in\mathbb R^{m\times k}$为
\begin{equation}
 (\bm A\bm B)_{i,j}=\sum_{\ell=1}^d\bm A_{i,\ell}\bm B_{\ell,j}.
 \label{eq:prelim-matrix-composition}
\end{equation}
左矩阵的列数必须等于右矩阵的行数；第$(i,j)$项取左矩阵的第$i$行与右矩阵的第$j$列，逐项相乘后相加。下面逐项计算一个二阶乘积。
\begin{example}[二阶矩阵的乘积]
\label{ex:ls-matrix-product-2x2}
给定：$\bm A=\left(\begin{smallmatrix}1&1\\0&1\end{smallmatrix}\right)$，$\bm B=\left(\begin{smallmatrix}1&0\\0&2\end{smallmatrix}\right)$。
计算：第$(i,j)$项是$\bm A$的第$i$行与$\bm B$的第$j$列的内积，
\[
 \bm A\bm B
 =\begin{pmatrix}1\cdot1+1\cdot0&1\cdot0+1\cdot2\\0\cdot1+1\cdot0&0\cdot0+1\cdot2\end{pmatrix}
 =\begin{pmatrix}1&2\\0&2\end{pmatrix}.
\]
结论：$\bm A\bm B=\left(\begin{smallmatrix}1&2\\0&2\end{smallmatrix}\right)$，形状为$2\times2$。
\end{example}
$\bm B$只有一列时，式~\eqref{eq:prelim-matrix-composition}就是式~\eqref{eq:prelim-matrix-prediction}。两个向量的内积$\bm x^\top\bm y$是$1\times d$与$d\times1$矩阵的乘积，结果是$1\times1$矩阵，我们把它与它唯一的元素等同看待。

这个定义来自两次加权求和的接连进行。

设$\bm x\in\mathbb R^k$，$\bm y=\bm B\bm x$、$\bm z=\bm A\bm y$，即先用$\bm B$再用$\bm A$，这是第~\ref{sec:prelim-map-operations}节意义下两个函数的复合。逐坐标代入并交换有限求和的次序，
\begin{equation}
 z_i=\sum_{\ell=1}^d a_{i,\ell}y_\ell
 =\sum_{\ell=1}^d a_{i,\ell}\sum_{j=1}^k b_{\ell,j}x_j
 =\sum_{j=1}^k\left(\sum_{\ell=1}^d a_{i,\ell}b_{\ell,j}\right)x_j,
 \label{eq:linear-algebra-matrix-product}
\end{equation}
括号中的数正是$(\bm A\bm B)_{i,j}$。所以$(\bm A\bm B)\bm x=\bm A(\bm B\bm x)$：直接用乘积作用，等于先用$\bm B$、再用$\bm A$作用。被求和的公共维数$d$就是被消去的中间结果$\bm y$的维数。

非方阵相乘时，写出每个输入和输出的形状尤其有用。以$2\times3$矩阵乘$3\times2$矩阵为例，
\[
 \begin{aligned}
 &\begin{pmatrix}1&2&0\\0&1&3\end{pmatrix}
 \begin{pmatrix}1&0\\2&1\\1&-1\end{pmatrix}\\
 &=\begin{pmatrix}1\cdot1+2\cdot2+0\cdot1&1\cdot0+2\cdot1+0\cdot(-1)\\
 0\cdot1+1\cdot2+3\cdot1&0\cdot0+1\cdot1+3\cdot(-1)\end{pmatrix}\\
 &=\begin{pmatrix}5&2\\5&-2\end{pmatrix}.
 \end{aligned}
\]
两个中间维数都是$3$，因此乘法有定义；输出保留左矩阵的$2$行和右矩阵的$2$列。若把其中一个$3$改为其他数，逐项配对便无法完成。

\paragraph{乘法的性质}
矩阵乘法保留结合律和分配律，但通常不满足交换律。下面逐项核对这些规则，同时检查矩阵形状是否相容。

\paragraph{单位乘}
\[
\bm I_m\bm A=\bm A=\bm A\bm I_d,\qquad
\bm A\in\mathbb R^{m\times d}.
\]
$\bm I_m$的第$i$行是$\bm e_i^\top$，它与$\bm A$的第$j$列相乘取出$a_{i,j}$；右乘同理。单位矩阵在矩阵乘法中的作用相当于数$1$在实数乘法中的作用。

\paragraph{零乘}
\[
 \begin{gathered}
\bm0\bm A=\bm0\\
\bm A\bm0=\bm0.
 \end{gathered}
\]
乘积的每个元素都是若干乘积之和，其中一个因子全为零时，每一项都为零。这里各个零矩阵的形状由相乘的另一个因子决定。

反过来不成立：$\bm A\bm B=\bm0$不能推出$\bm A=\bm0$或$\bm B=\bm0$，下面是一个反例。
\begin{example}[反例：两个非零矩阵的乘积为零]
\label{ex:ls-zero-product-nonzero-factors}
给定：$\bm A=\left[\begin{smallmatrix}1&0\\0&0\end{smallmatrix}\right]$，$\bm B=\left[\begin{smallmatrix}0&0\\1&0\end{smallmatrix}\right]$，两者都不是零矩阵。
计算：
\[
 \bm A\bm B=\begin{bmatrix}1\cdot0+0\cdot1&1\cdot0+0\cdot0\\0\cdot0+0\cdot1&0\cdot0+0\cdot0\end{bmatrix}=\begin{bmatrix}0&0\\0&0\end{bmatrix}.
\]
结论：$\bm A\bm B=\bm0$，而$\bm A\neq\bm0$、$\bm B\neq\bm0$。
\end{example}
原因可以按列看出。由乘积的元素公式，$\bm A\bm B$的第$j$列就是$\bm A$乘$\bm B$的第$j$列，所以$\bm A\bm B=\bm0$恰好意味着$\bm B$的每一列$\bm b$都满足$\bm A\bm b=\bm0$：右因子产生的结果被左因子完全消去。上例中$\bm A$只保留第一个坐标、把第二个坐标清零，而$\bm B$的第一列$(0,1)^\top$只有第二个坐标非零，第二列为零。

\paragraph{结合律}
\[
 (\bm A\bm B)\bm C=\bm A(\bm B\bm C).
\]
按式~\eqref{eq:prelim-matrix-composition}展开，两边的第$(i,j)$项都是$\sum_k\sum_l a_{i,k}b_{k,l}c_{l,j}$，只是两个有限求和的次序不同。这也是函数复合满足结合律（第~\ref{sec:prelim-map-operations}节）在矩阵上的翻译。由结合律，多个相容矩阵连乘时不必标出括号。

\paragraph{分配律与数乘相容}
\[
 \begin{gathered}
\bm A(\bm B+\bm C)=\bm A\bm B+\bm A\bm C\\
(\bm A+\bm B)\bm C=\bm A\bm C+\bm B\bm C\\
c(\bm A\bm B)=(c\bm A)\bm B=\bm A(c\bm B),
 \end{gathered}
\]
逐项展开有限和即得。把右因子取为列向量，得到矩阵与向量的乘积对输入保持线性组合：
\begin{equation}
\bm A(a\bm u+b\bm v)=a\bm A\bm u+b\bm A\bm v,\qquad
\bm u,\bm v\in\mathbb R^d,\quad a,b\in\mathbb R.
\label{eq:prelim-matrix-linearity}
\end{equation}

由这一性质，函数$\bm u\mapsto\bm A\bm u$由它在$\bm e_1,\ldots,\bm e_d$上的输出完全确定：因为$\bm u=\sum_ju_j\bm e_j$，所以$\bm A\bm u=\sum_ju_j\bm A\bm e_j$，而这些输出$\bm A\bm e_j$就是$\bm A$的各列。第~\ref{chap:linear-representation}章把保持线性组合的函数称为线性映射，并说明每个这样的映射都能由一张矩阵记录。

\paragraph{分块乘法}
把矩阵按行、列分成尺寸相容的块，乘法规则对块照样成立：
\[
\begin{gathered}
\begin{bmatrix}\bm A&\bm B\end{bmatrix}\begin{bmatrix}\bm C\\\bm D\end{bmatrix}=\bm A\bm C+\bm B\bm D\\
\begin{bmatrix}\bm A\\\bm B\end{bmatrix}\bm C=\begin{bmatrix}\bm A\bm C\\\bm B\bm C\end{bmatrix},
\end{gathered}
\]
\[
 \bm A\begin{bmatrix}\bm C&\bm D\end{bmatrix}=\begin{bmatrix}\bm A\bm C&\bm A\bm D\end{bmatrix},
\]
以及$2\times2$分块的乘积
\[
 \begin{bmatrix}\bm A_{1,1}&\bm A_{1,2}\\\bm A_{2,1}&\bm A_{2,2}\end{bmatrix}
 \begin{bmatrix}\bm B_{1,1}&\bm B_{1,2}\\\bm B_{2,1}&\bm B_{2,2}\end{bmatrix}
 =\begin{bmatrix}
  \bm A_{1,1}\bm B_{1,1}+\bm A_{1,2}\bm B_{2,1}&\bm A_{1,1}\bm B_{1,2}+\bm A_{1,2}\bm B_{2,2}\\
  \bm A_{2,1}\bm B_{1,1}+\bm A_{2,2}\bm B_{2,1}&\bm A_{2,1}\bm B_{1,2}+\bm A_{2,2}\bm B_{2,2}
 \end{bmatrix}.
\]
条件是每个相乘的块对中，左块的列数等于右块的行数。第一式成立，是因为式~\eqref{eq:prelim-matrix-composition}对中间索引$\ell$的求和可以分段进行：前一段$\ell$对应$\bm A$的列与$\bm C$的行，后一段对应$\bm B$的列与$\bm D$的行，两段之和逐项等于整段求和。后两式只是把结果的行或列分组，$2\times2$分块由这三式合成。块可以像数一样参与运算，只是块之间的次序不能交换。

把块分到最细，得到三种常用的特例。右因子按列分块，得$(\bm A\bm B)_{:,j}=\bm A\bm B_{:,j}$：$\bm A$分别乘$\bm B$的每一列。左因子按行分块，得$(\bm A\bm B)_{i,:}=\bm A_{i,:}\bm B$：结果的第$i$行是$\bm B$各行以$\bm A_{i,:}$的元素为系数的线性组合。

第三种特例把中间维数$d$分成长度为$1$的$d$段，这时第一式给出
\begin{equation}
 \bm A\bm B=\sum_{\ell=1}^d\bm A_{:,\ell}\bm B_{\ell,:},
 \label{eq:ls-product-outer-sum}
\end{equation}
其中一列$\bm A_{:,\ell}\in\mathbb R^{m\times1}$乘一行$\bm B_{\ell,:}\in\mathbb R^{1\times k}$，得到$(i,j)$元素为$a_{i,\ell}b_{\ell,j}$的$m\times k$矩阵。式~\eqref{eq:ls-product-outer-sum}是后面低秩近似的出发点。下面用示例~\ref{ex:ls-matrix-product-2x2}的乘积核对这个公式。
\begin{example}[把乘积写成列乘行之和]
\label{ex:ls-product-as-outer-sum}
给定：示例~\ref{ex:ls-matrix-product-2x2}中的$\bm A=\left(\begin{smallmatrix}1&1\\0&1\end{smallmatrix}\right)$与$\bm B=\left(\begin{smallmatrix}1&0\\0&2\end{smallmatrix}\right)$。
计算：$\bm A$的两列为$(1,0)^\top$与$(1,1)^\top$，$\bm B$的两行为$(1,0)$与$(0,2)$，
\[
 \bm A\bm B=\begin{pmatrix}1\\0\end{pmatrix}\begin{pmatrix}1&0\end{pmatrix}+\begin{pmatrix}1\\1\end{pmatrix}\begin{pmatrix}0&2\end{pmatrix}
 =\begin{pmatrix}1&0\\0&0\end{pmatrix}+\begin{pmatrix}0&2\\0&2\end{pmatrix}=\begin{pmatrix}1&2\\0&2\end{pmatrix}.
\]
结论：两个“列乘行”矩阵之和与逐项计算的乘积相同。
\end{example}

\paragraph{交换律不成立}
一般$\bm A\bm B\neq\bm B\bm A$。形状不同时，两个乘积甚至不一定都有定义；即使两个因子是同阶方阵，两种次序的乘积也可能不同，下面是一个反例。
\begin{example}[反例：先交换再保留与先保留再交换]
\label{ex:ls-noncommuting-pair}
给定：只保留第一个坐标的$\bm P=\left[\begin{smallmatrix}1&0\\0&0\end{smallmatrix}\right]$，交换两个坐标的$\bm S=\left[\begin{smallmatrix}0&1\\1&0\end{smallmatrix}\right]$。
计算：
\[
\begin{gathered}
\bm P\bm S=\begin{bmatrix}0&1\\0&0\end{bmatrix}\\
\bm S\bm P=\begin{bmatrix}0&0\\1&0\end{bmatrix}.
\end{gathered}
\]
结论：$\bm P\bm S\neq\bm S\bm P$。作用在$\bm u\in\mathbb R^2$上，$\bm P\bm S\bm u=(u_2,0)^\top$是先交换、再保留第一个坐标，$\bm S\bm P\bm u=(0,u_1)^\top$是先保留、再交换，两种次序的结果不同。
\end{example}
书写$\bm A\bm B\bm x$时，作用顺序是从右向左：先$\bm B$，后$\bm A$。

\paragraph{消去律不成立}
$\bm A\bm B=\bm A\bm C$且$\bm A\neq\bm0$不能推出$\bm B=\bm C$。原因是$\bm A\bm B=\bm A\bm C$等价于$\bm A(\bm B-\bm C)=\bm0$，而乘积为零不能推出$\bm B-\bm C=\bm0$。下面用示例~\ref{ex:ls-zero-product-nonzero-factors}构造反例。
\begin{example}[反例：非零的左因子不能消去]
\label{ex:ls-cancellation-fails}
给定：示例~\ref{ex:ls-zero-product-nonzero-factors}中的$\bm A=\left[\begin{smallmatrix}1&0\\0&0\end{smallmatrix}\right]$与$\bm B=\left[\begin{smallmatrix}0&0\\1&0\end{smallmatrix}\right]$，以及$\bm C=\bm0$。
验证：$\bm A\bm B=\bm0=\bm A\bm C$，而$\bm B\neq\bm C$。
结论：虽然$\bm A\bm B=\bm A\bm C$且$\bm A\neq\bm0$，却不能消去$\bm A$得到$\bm B=\bm C$。
\end{example}
若$\bm A$是可逆方阵，在$\bm A\bm B=\bm A\bm C$两边左乘$\bm A^{-1}$即得$\bm B=\bm C$。所以左因子可逆时可以消去，一般的非零因子则不能。

常见矩阵参与乘法时，结果有固定模式。下面逐类给出。

\paragraph{对角矩阵的乘法}
对$\bm D=\operatorname{diag}(d_1,\ldots,d_n)$，
\[
 \begin{gathered}
(\bm D\bm A)_{i,j}=d_i\bm A_{i,j}\\
(\bm A\bm D)_{i,j}=\bm A_{i,j}d_j,
 \end{gathered}
\]
\[
 \operatorname{diag}(d_1,\ldots,d_n)\operatorname{diag}(e_1,\ldots,e_n)=\operatorname{diag}(d_1e_1,\ldots,d_ne_n).
\]
左乘对角矩阵缩放各行，右乘缩放各列：乘积公式中只有$\ell=i$（左乘时）或$\ell=j$（右乘时）的一项留下。作用在向量上，$\bm D\bm u=(d_1u_1,\ldots,d_nu_n)^\top$是逐坐标缩放。两个对角矩阵相乘时对角元逐个相乘，所以对角矩阵之间可以交换次序。下面的例子比较左乘与右乘。
\begin{example}[左乘与右乘对角矩阵]
\label{ex:ls-diagonal-scaling}
给定：$\bm D=\operatorname{diag}(2,3)$，$\bm A=\left[\begin{smallmatrix}1&2\\3&4\end{smallmatrix}\right]$。
计算：
\[
 \begin{aligned}
 \bm D\bm A&=\begin{bmatrix}2\cdot1&2\cdot2\\3\cdot3&3\cdot4\end{bmatrix}=\begin{bmatrix}2&4\\9&12\end{bmatrix},\\
 \bm A\bm D&=\begin{bmatrix}1\cdot2&2\cdot3\\3\cdot2&4\cdot3\end{bmatrix}=\begin{bmatrix}2&6\\6&12\end{bmatrix}.
 \end{aligned}
\]
结论：左乘把第一、二行分别乘以$2$、$3$，右乘把第一、二列分别乘以$2$、$3$。两个结果不同，所以对角矩阵与一般矩阵相乘时仍要区分左右。
\end{example}

\paragraph{三角矩阵的乘法}
若$\bm A,\bm B$都是$n$阶上三角矩阵，则$\bm A\bm B$也是上三角矩阵，并且
\[
 (\bm A\bm B)_{i,i}=\bm A_{i,i}\bm B_{i,i}.
\]
在$(\bm A\bm B)_{i,j}=\sum_\ell\bm A_{i,\ell}\bm B_{\ell,j}$中，$\bm A_{i,\ell}$只在$\ell\geq i$时可能非零，$\bm B_{\ell,j}$只在$\ell\leq j$时可能非零，所以只有$i\leq\ell\leq j$的项可能非零：$i>j$时没有这样的$\ell$，乘积元素为零；$i=j$时只剩$\ell=i$一项。下三角矩阵同理：两个下三角矩阵的乘积仍是下三角矩阵，对角元逐个相乘。下面是一个二阶算例。
\begin{example}[两个上三角矩阵的乘积]
\label{ex:ls-upper-triangular-product}
给定：$\bm A=\left[\begin{smallmatrix}1&2\\0&3\end{smallmatrix}\right]$，$\bm B=\left[\begin{smallmatrix}4&5\\0&6\end{smallmatrix}\right]$。
计算：
\[
 \bm A\bm B=\begin{bmatrix}1\cdot4+2\cdot0&1\cdot5+2\cdot6\\0\cdot4+3\cdot0&0\cdot5+3\cdot6\end{bmatrix}=\begin{bmatrix}4&17\\0&18\end{bmatrix}.
\]
结论：乘积仍是上三角矩阵，对角元$4=1\cdot4$、$18=3\cdot6$是两个因子对角元的乘积；右上角的$17$则同时用到了两个因子的非对角元。
\end{example}

\paragraph{置换矩阵的乘法}
设$\bm P=\left[\begin{smallmatrix}\bm e_{\sigma(1)}&\cdots&\bm e_{\sigma(n)}\end{smallmatrix}\right]$，其中$\sigma$是$\{1,\ldots,n\}$上的排列，则
\[
 \begin{gathered}
(\bm P\bm A)_{\sigma(i),:}=\bm A_{i,:}\\
(\bm A\bm P)_{:,j}=\bm A_{:,\sigma(j)}.
 \end{gathered}
\]
即左乘把$\bm A$的第$i$行移到第$\sigma(i)$行，右乘使$\bm A\bm P$的第$j$列取自$\bm A$的第$\sigma(j)$列。例如$\left[\begin{smallmatrix}0&1\\1&0\end{smallmatrix}\right]$左乘交换两行、右乘交换两列。

两个置换矩阵的乘积仍是置换矩阵：设$\bm P'$是另一个置换矩阵，$\bm P'\bm P$的第$j$列是$\bm P'\bm e_{\sigma(j)}$，即$\bm P'$的第$\sigma(j)$列；它是标准坐标向量，而且不同的$j$取到$\bm P'$的不同列，所以这些列互不相同。

\paragraph{对称矩阵的乘积}
若$\bm A,\bm B$是同阶对称矩阵，则
\[
 (\bm A\bm B)^\top=\bm B^\top\bm A^\top=\bm B\bm A,
\]
所以$\bm A\bm B$对称当且仅当$\bm A\bm B=\bm B\bm A$。对称矩阵的乘积一般不对称，下面是一个反例。
\begin{example}[反例：两个对称矩阵的乘积不对称]
\label{ex:ls-symmetric-product-not-symmetric}
给定：对称矩阵$\bm A=\left[\begin{smallmatrix}1&1\\1&0\end{smallmatrix}\right]$与$\bm B=\operatorname{diag}(1,2)$。
计算：右乘对角矩阵缩放各列，
\[
 \bm A\bm B=\begin{bmatrix}1\cdot1&1\cdot2\\1\cdot1&0\cdot2\end{bmatrix}=\begin{bmatrix}1&2\\1&0\end{bmatrix}.
\]
结论：$(\bm A\bm B)_{1,2}=2\neq1=(\bm A\bm B)_{2,1}$，乘积不对称。这与上面的判据一致：$\bm B\bm A=\left[\begin{smallmatrix}1&1\\2&0\end{smallmatrix}\right]\neq\bm A\bm B$。
\end{example}


\subsubsection{逐项乘积}
\label{sec:ls-hadamard-product}

矩阵乘法汇总共同索引；另一种常用的“乘法”不汇总任何索引，只把两个同形数组在对应位置配对。

设$\bm A,\bm B\in\mathbb R^{m\times n}$，\term{逐元素乘积}（Hadamard product）记作$\bm A\odot\bm B$，
\begin{equation}
 (\bm A\odot\bm B)_{i,j}=a_{i,j}b_{i,j},
 \label{eq:ls-hadamard}
\end{equation}
输出仍是$m\times n$。同维向量和更高阶的同形数组也按同样的方式逐项相乘。

同样两个矩阵，用两种乘法得到的结果完全不同，下面的例子说明这一点。
\begin{example}[逐项乘积与矩阵乘积的对比]
\label{ex:ls-hadamard-vs-matrix-product}
给定：$\bm A=\left[\begin{smallmatrix}1&2\\3&4\end{smallmatrix}\right]$，$\bm B=\left[\begin{smallmatrix}0&1\\1&0\end{smallmatrix}\right]$。
计算：
\[
 \begin{aligned}
 \bm A\odot\bm B&=\begin{bmatrix}1\cdot0&2\cdot1\\3\cdot1&4\cdot0\end{bmatrix}=\begin{bmatrix}0&2\\3&0\end{bmatrix},\\
 \bm A\bm B&=\begin{bmatrix}1\cdot0+2\cdot1&1\cdot1+2\cdot0\\3\cdot0+4\cdot1&3\cdot1+4\cdot0\end{bmatrix}=\begin{bmatrix}2&1\\4&3\end{bmatrix}.
 \end{aligned}
\]
结论：逐项乘积把$\bm B$为零的位置置零，矩阵乘积则交换了$\bm A$的两列；两个结果在每个位置上都不同。
\end{example}
固定掩码$\bm M$后，$\bm X\mapsto\bm M\odot\bm X$把$\bm M$为零的位置置零、把其余位置按系数缩放，适合表达“挑选或调制另一个数组”。

\paragraph{逐项乘积的性质}

逐项乘积的规则可以逐个位置从实数运算推出。以下矩阵的形状都相同。

\paragraph{交换律}
\[
 \bm A\odot\bm B=\bm B\odot\bm A.
\]
式~\eqref{eq:ls-hadamard}右端只是实数乘法，按位置逐项验证即得。逐项乘积可交换而矩阵乘积不可交换，再次说明两者是不同的运算。

\paragraph{结合律}
\[
 (\bm A\odot\bm B)\odot\bm C=\bm A\odot(\bm B\odot\bm C),
\]
两边的$(i,j)$元素都是$a_{i,j}b_{i,j}c_{i,j}$。

\paragraph{分配律}
\[
 \bm A\odot(\bm B+\bm C)=\bm A\odot\bm B+\bm A\odot\bm C,
\]
两边的$(i,j)$元素都是$a_{i,j}(b_{i,j}+c_{i,j})$。

\paragraph{单位元是全$1$数组}
全部元素为$1$的同形数组$\bm J$满足$\bm J\odot\bm A=\bm A$。单位矩阵不是逐项乘积的单位元：$\bm I_2\odot\bm A$只保留$\bm A$的对角元，把非对角元都变为零。

\paragraph{对每个因子分别线性}
固定一个因子时，$\bm X\mapsto\bm M\odot\bm X$保持线性组合。两个因子都作为未知量时，乘积只对各自分别线性：同时把两者放大两倍，结果放大四倍。第~\ref{chap:linear-representation}章换一组记录方向以后会看到，矩阵乘积对应的复合关系保持不变，逐项乘积则随记录方式改变。


\subsection{初等行变换与阶梯形}
\label{sec:ls-elementary-row-operations}

求逆时，我们需要把一个矩阵逐步化成容易辨认的形式。每一步都应当可以撤销，这样才不会丢失原矩阵中的信息。对矩阵的行允许三种\term{初等行变换}（elementary row operation）：交换两行；将一行乘以非零数；将一行加上另一行的倍数。它们的逆操作依次是再次交换同两行、乘以原系数的倒数，以及减去原来加上的倍数。

每种操作都可写成左乘一个初等矩阵：从单位矩阵出发，对其执行同样的行变换，所得矩阵记为$\bm E$，则$\bm E\bm A$恰好对$\bm A$执行该变换。例如交换两行的矩阵交换乘积的对应两行。逆操作也对应初等矩阵，因此$\bm E$可逆。

化简时从左到右寻找\term{主元}（pivot）。先在当前列尚未使用的行中找一个非零元素；若有必要，交换行把它移到当前行，再用它消去同列下方元素。若该列在余下各行全为零，就跳过这一列并继续向右。重复到没有可用行或列为止，得到的\term{行阶梯形}（row echelon form）中，零行位于底部，每个非零行的首个非零元位于上一行首个非零元的右侧。再把每个主元化为$1$，消去其上方元素，就得到\term{行最简形}（reduced row echelon form）。

处理方程$\bm A\bm x=\bm b$时，把右端项作为最后一列附在$\bm A$后面，得到增广矩阵$[\bm A\mid\bm b]$，并对整行执行同一操作。交换等式、将等式乘以非零数、把一个等式的倍数加到另一个等式，都不改变解；逆操作保证新的等式组也能还原为原等式组。下面在一个零候选主元和一列冗余的例子中逐步计算。
\begin{example}[换行、消元与自由变量]
\label{ex:ls-elementary-row-reduction}
考虑以下增广矩阵。第一列的首个元素为零，先交换前两行，再消去第三行的第一列：
\[
 \left[\begin{array}{ccc|c}0&1&2&3\\1&2&4&6\\2&4&8&12\end{array}\right]
 \longrightarrow
 \left[\begin{array}{ccc|c}1&2&4&6\\0&1&2&3\\2&4&8&12\end{array}\right]
 \longrightarrow
 \left[\begin{array}{ccc|c}1&2&4&6\\0&1&2&3\\0&0&0&0\end{array}\right].
\]
这已是阶梯形。将第一行减去第二行的两倍，得到行最简形
\[
 \left[\begin{array}{ccc|c}1&0&0&0\\0&1&2&3\\0&0&0&0\end{array}\right].
\]
第一、二列是主元列，第三列没有主元。令自由变量$x_3=t$，便读出$x_1=0$、$x_2=3-2t$；任意实数$t$都给出原方程组的一个解。最后的零行表示$0=0$，不会增加约束；若出现系数全为零而右端非零的行，则方程组无解。
\end{example}

阶梯形使主元和自由变量可见，但主元的选择会影响实际数值计算的误差。第~\ref{sec:ls-direct-methods}节会在求解方程时讨论选主元和计算代价。

\subsection{逆矩阵}
\label{sec:prelim-inverse-matrices}

对$n$阶方阵$\bm A$，若存在$n$阶方阵$\bm B$满足
\[
 \bm A\bm B=\bm B\bm A=\bm I_n,
\]
就称$\bm A$可逆，$\bm B$为\term{逆矩阵}（inverse matrix），记作$\bm A^{-1}$。逆矩阵撤销乘$\bm A$的效果。$\bm A$可逆时，若$\bm A\bm x=\bm y$，两边左乘$\bm A^{-1}$得$\bm x=\bm A^{-1}\bm y$，所以方程至多有一个解；反过来，$\bm x=\bm A^{-1}\bm y$满足$\bm A\bm x=\bm A\bm A^{-1}\bm y=\bm y$。所以方程$\bm A\bm x=\bm y$有唯一解$\bm x=\bm A^{-1}\bm y$。

逆矩阵若存在则唯一：若$\bm B,\bm C$都满足条件，则$\bm B=\bm B(\bm A\bm C)=(\bm B\bm A)\bm C=\bm C$，这里用到前面建立的结合律。因此记号$\bm A^{-1}$不会有歧义。

怎样求逆取决于矩阵的结构。下面从最简单的对角矩阵开始，依次给出二阶方阵的公式与一般方阵的做法。

\paragraph{对角矩阵的逆}
每个$d_i\neq0$时逐个取倒数：
\[
 \operatorname{diag}(d_1,\ldots,d_n)^{-1}=\operatorname{diag}(1/d_1,\ldots,1/d_n).
\]
两个对角矩阵相乘时对角元逐个相乘（前面已验证），所以两侧乘积都是$\bm I_n$。例如$\operatorname{diag}(2,1)^{-1}=\operatorname{diag}(1/2,1)$。若某个$d_i=0$，左乘$\bm D=\operatorname{diag}(d_1,\ldots,d_n)$把任何$\bm B$的第$i$行乘以$d_i$，所以$\bm D\bm B$的第$i$行是$\bm0$，不可能等于$\bm I_n$的第$i$行；这时$\bm D$不可逆。

\paragraph{二阶方阵的逆}
$ad-bc\neq0$时，
\[
 \begin{pmatrix}a&b\\c&d\end{pmatrix}^{-1}=\frac1{ad-bc}\begin{pmatrix}d&-b\\-c&a\end{pmatrix}.
\]
直接相乘即可验证两侧乘积都是$\bm I_2$：例如左上元素为$(ad-bc)/(ad-bc)=1$，右上元素为$(-ab+ba)/(ad-bc)=0$。分母$ad-bc$正是后面的二阶行列式。

\paragraph{一般方阵的逆}
一般方阵没有这样简短的公式，可以通过解方程求逆：对任意$\bm y$解$\bm A\bm x=\bm y$，把$\bm x$写成$\bm y$各坐标的线性组合，系数排成的矩阵就是$\bm A^{-1}$。下面用一个二阶例子演示，并与二阶公式对照。
\begin{example}[解方程求逆]
\label{ex:ls-inverse-by-solving}
给定：$\bm A=\left[\begin{smallmatrix}2&1\\0&3\end{smallmatrix}\right]$。
计算：方程$\bm A\bm x=\bm y$即$y_1=2x_1+x_2$、$y_2=3x_2$。由第二式得$x_2=y_2/3$，代入第一式得$x_1=(y_1-x_2)/2=y_1/2-y_2/6$。把系数排成矩阵，
\[
 \bm A^{-1}=\begin{bmatrix}1/2&-1/6\\0&1/3\end{bmatrix}.
\]
用二阶公式核对：$ad-bc=2\cdot3-1\cdot0=6$，公式给出$\frac16\left[\begin{smallmatrix}3&-1\\0&2\end{smallmatrix}\right]$，与上式相同。
结论：右上角的$-1/6$不是原元素$1$的倒数。一般矩阵把各个坐标掺在一起，逆矩阵要作为整体求出，不能逐个元素取倒数。
\end{example}
上述求法也可以用初等行变换同时处理所有右端项。把$\bm A$与单位矩阵并排，得到$[\bm A\mid\bm I_n]$，对整个增广矩阵作行变换。若左侧最终化为$\bm I_n$，这些操作合起来是某个可逆矩阵$\bm E$的左乘，于是$\bm E\bm A=\bm I_n$、右侧为$\bm E$，故$\bm E=\bm A^{-1}$。例如
\[
 \begin{aligned}
 &\left[\begin{array}{cc|cc}2&1&1&0\\0&3&0&1\end{array}\right]\\
 &\longrightarrow
 \left[\begin{array}{cc|cc}1&1/2&1/2&0\\0&1&0&1/3\end{array}\right]\\
 &\longrightarrow
 \left[\begin{array}{cc|cc}1&0&1/2&-1/6\\0&1&0&1/3\end{array}\right],
 \end{aligned}
\]
右半部分与前面的逆矩阵相同。如果方阵化简后不能在每列找到主元，左侧出现零行；由于所有行变换都可逆，这样的矩阵不能有逆。计算方程的解时，通常直接对右端项消元，不必先求出整个逆矩阵。哪些方阵可逆，第~\ref{sec:ls-injective-surjective-invertible}节用零空间与秩回答，第~\ref{sec:ls-determinant-invertibility}节再给出行列式判据。

\paragraph{单侧逆}
矩形矩阵不能有双侧的逆，但可能只从一侧抵消。对$\bm A\in\mathbb R^{m\times d}$，满足$\bm L\bm A=\bm I_d$的$\bm L\in\mathbb R^{d\times m}$称为左逆，满足$\bm A\bm R=\bm I_m$的$\bm R\in\mathbb R^{d\times m}$称为右逆。有左逆时，若$\bm A\bm x=\bm y$有解，两边左乘$\bm L$得$\bm x=\bm L\bm y$，所以解只能是$\bm L\bm y$。有右逆时，$\bm A(\bm R\bm y)=\bm I_m\bm y=\bm y$，所以$\bm x=\bm R\bm y$总是一个解。

与方阵的逆不同，单侧逆一般不唯一，下面的例子说明这一点。
\begin{example}[反例：左逆不唯一]
\label{ex:ls-left-inverse-not-unique}
给定：$\bm A=(1,0)^\top\in\mathbb R^{2\times1}$，以及任意实数$c$。
验证：取$\bm L=(1,c)\in\mathbb R^{1\times2}$，则$\bm L\bm A=1\cdot1+c\cdot0=1=\bm I_1$。
结论：每个$c$都给出$\bm A$的一个左逆，所以$\bm A$的左逆有无穷多个。
\end{example}
矩形矩阵何时有左逆或右逆，第~\ref{chap:linear-representation}章用核与像回答；第~\ref{chap:linear-equations-and-optimization}章讨论怎样实际求解，以及多解时选哪一个。

设下面的$\bm A,\bm B$都是$n$阶可逆方阵。

\paragraph{乘积的逆}
\begin{equation}
 (\bm A\bm B)^{-1}=\bm B^{-1}\bm A^{-1}.
 \label{eq:ls-inverse-product}
\end{equation}
先作用的后撤销，就像先穿袜子再穿鞋，脱的时候要先脱鞋。验证：$(\bm A\bm B)(\bm B^{-1}\bm A^{-1})=\bm A(\bm B\bm B^{-1})\bm A^{-1}=\bm A\bm A^{-1}=\bm I_n$，反方向$(\bm B^{-1}\bm A^{-1})(\bm A\bm B)=\bm I_n$同理。

\paragraph{逆的逆}
\[
 (\bm A^{-1})^{-1}=\bm A.
\]
定义中$\bm A$与$\bm A^{-1}$的地位对称：$\bm A^{-1}\bm A=\bm A\bm A^{-1}=\bm I_n$同样说明$\bm A$是$\bm A^{-1}$的逆。

\paragraph{数乘的逆}
\[
 (c\bm A)^{-1}=c^{-1}\bm A^{-1}\qquad(c\neq0).
\]
直接相乘：$(c\bm A)(c^{-1}\bm A^{-1})=(cc^{-1})\bm A\bm A^{-1}=\bm I_n$，反方向同理。

\paragraph{转置的逆}
\[
 (\bm A^\top)^{-1}=(\bm A^{-1})^\top.
\]
由乘积的转置，$\bm A^\top(\bm A^{-1})^\top=(\bm A^{-1}\bm A)^\top=\bm I^\top=\bm I$，反方向同理，所以$(\bm A^{-1})^\top$就是$\bm A^\top$的逆。

\subsection{矩阵的幂}
\label{sec:ls-matrix-power-rules}

定常迭代$\bm x^{(k+1)}=\bm A\bm x^{(k)}$每一步都乘同一个$\bm A$。对$d$阶方阵，$\bm A^k$定义为$k$个$\bm A$连乘（$k\geq1$），并规定$\bm A^0=\bm I_d$；$\bm A$可逆时再定义负幂$\bm A^{-k}=(\bm A^{-1})^k$。于是迭代的第$k$步是$\bm x^{(k)}=\bm A^k\bm x^{(0)}$。非方阵不能与自身相乘，所以没有矩阵幂。

幂是反复的矩阵乘法，不是逐元素乘方，下面的例子说明两者的差别。
\begin{example}[矩阵的平方不是逐元素平方]
\label{ex:ls-matrix-square-shear}
给定：$\bm A=\left(\begin{smallmatrix}1&1\\0&1\end{smallmatrix}\right)$。
计算：
\[
 \bm A^2=\begin{pmatrix}1&1\\0&1\end{pmatrix}\begin{pmatrix}1&1\\0&1\end{pmatrix}
 =\begin{pmatrix}1\cdot1+1\cdot0&1\cdot1+1\cdot1\\0\cdot1+1\cdot0&0\cdot1+1\cdot1\end{pmatrix}
 =\begin{pmatrix}1&2\\0&1\end{pmatrix}.
\]
结论：右上角的元素变为$2$，而逐元素平方会得到$1^2=1$；继续相乘，右上角每次增加$1$。
\end{example}

同一矩阵的整数幂按结合律满足$\bm A^r\bm A^s=\bm A^{r+s}$和$(\bm A^r)^s=\bm A^{rs}$；出现负指数时要求$\bm A$可逆。但不同矩阵不一定交换，因此
\[
 (\bm A+\bm B)^2=\bm A^2+\bm A\bm B+\bm B\bm A+\bm B^2,
\]
只有$\bm A\bm B=\bm B\bm A$时才能将中间两项合成$2\bm A\bm B$。这与标量平方公式的差别来自乘法次序。

两类幂关系经常出现。满足$\bm A^2=\bm A$的方阵称为\term{幂等}（idempotent）矩阵，乘一次与乘多次效果相同。存在正整数$s$使$\bm N^s=\bm0$的方阵称为\term{幂零}（nilpotent）矩阵。下面各举一例。
\begin{example}[幂等矩阵与幂零矩阵]
\label{ex:ls-idempotent-nilpotent}
给定：$\bm P=\operatorname{diag}(1,0)$，$\bm N=\left(\begin{smallmatrix}0&1\\0&0\end{smallmatrix}\right)$。
计算：对$\bm u\in\mathbb R^2$，$\bm P\bm u=(u_1,0)^\top$只保留第一个坐标，再作用一次结果不变，所以$\bm P^2=\bm P$。$\bm N\bm u=(u_2,0)^\top$把第二个坐标移到第一个位置、把第二个位置清零；再作用一次得$\bm N(u_2,0)^\top=(0,0)^\top$。直接相乘也得到
\[
 \bm N^2=\begin{pmatrix}0&1\\0&0\end{pmatrix}\begin{pmatrix}0&1\\0&0\end{pmatrix}=\begin{pmatrix}0&0\\0&0\end{pmatrix}.
\]
结论：$\bm P$是幂等矩阵；$\bm N\neq\bm0$，却满足$\bm N^2=\bm0$，是幂零矩阵。
\end{example}
$\bm A^k$在$k$很大时是什么样子，很难从元素看出，第~\ref{chap:linear-analysis}章借助特征向量回答。

\paragraph{幂的性质}

设$\bm A$为方阵，$m,n$为非负整数；$\bm A$可逆时，$m,n$可取任意整数。

\paragraph{同底幂相乘}
\[
 \bm A^m\bm A^n=\bm A^{m+n}.
\]
由结合律，两边都是同样个数的$\bm A$（或$\bm A^{-1}$）连乘；指数异号时，相邻的$\bm A\bm A^{-1}$逐对消去。交换$m,n$的角色可见，同一矩阵的不同次幂彼此可以交换。

\paragraph{幂的幂}
\[
 (\bm A^m)^n=\bm A^{mn}.
\]
左边是$n$组、每组$m$个因子的连乘，由结合律去掉括号即得。

\paragraph{负幂与逆相容}
\[
 (\bm A^k)^{-1}=(\bm A^{-1})^k=\bm A^{-k}.
\]
对$k$个相同因子反复使用式~\eqref{eq:ls-inverse-product}即得。
